\SourceLeaf{375}%
then \(DC=CL\) would have to hold, in contradiction with the result found before. We believe we have indicated what matters at this stage of the proof, as the author sees it, a little more clearly and precisely than he did himself. This, it seems to us, also makes the weakness of the proof easier to see. The assumption made here is clearly entirely arbitrary. In all isosceles triangles with angle \(A\) at the vertex and larger angles at the base, when the construction indicated in the theorem is carried out, the succession of the segments cut off, with respect to their staying equal, becoming larger or becoming smaller, must always, regardless of the size of the sides. necessarily be the same. That assumption cannot possibly be regarded as self-evident. Since the attempted proof that the third case (as well as the fourth and fifth) is impossible rests on this alone, and the entire article offers no other means of proving the omitted case impossible, we believe we have thus sufficiently justified the judgment stated above. We gladly give full credit, however, to the whole ingenious treatment in the articles that follow.

